% =====================================================================
% Consulta de avance - Agitador magnético con motor BLDC sensorless
% Compilar desde esta carpeta (dos pasadas para el total de slides):
%   pdflatex presentacion_agitador.tex
% Las figuras leen datos/aire.dat y datos/agua.dat. Para regenerarlos desde
% los CSV de la página /curvas (desde la raíz del repo):
%   python docs/presentacion/preparar_datos.py
% =====================================================================
\documentclass[aspectratio=169,10pt,t]{beamer}

\usepackage[utf8]{inputenc}
\usepackage[T1]{fontenc}
\usepackage[spanish,es-noshorthands]{babel}
\usepackage{lmodern}
\usepackage{amsmath,amssymb}
\usepackage{tikz}
\usetikzlibrary{arrows.meta,calc}
\usepackage{pgfplots}
\pgfplotsset{compat=1.18}

% ---------------------------------------------------------------------
% Estilo
% ---------------------------------------------------------------------
\definecolor{azul}{HTML}{1F3A5F}
\definecolor{acento}{HTML}{0F7C80}
\definecolor{rojo}{HTML}{B23A48}
\definecolor{gris}{HTML}{5F6B7A}
\definecolor{grisclaro}{HTML}{E6EAEE}
\definecolor{verde}{HTML}{2E7D4F}

\usefonttheme{professionalfonts}
\setbeamertemplate{navigation symbols}{}
\setbeamersize{text margin left=0.75cm, text margin right=0.75cm}
\setbeamercolor{normal text}{fg=black!85}
\setbeamercolor{structure}{fg=azul}
\setbeamercolor{alerted text}{fg=rojo}
\setbeamerfont{frametitle}{series=\bfseries, size=\large}
\setbeamertemplate{frametitle}{%
  \vspace*{0.3cm}%
  {\usebeamerfont{frametitle}\color{azul}\insertframetitle\par}%
  \nointerlineskip\vspace{4pt}%
  \hbox{\color{acento}\rule{1.6cm}{1.5pt}}%
  \vspace{0.1cm}}
\setbeamertemplate{itemize item}{\raisebox{0.15ex}{\color{acento}\scriptsize$\blacktriangleright$}}
\setbeamertemplate{enumerate item}{\color{acento}\bfseries\insertenumlabel.}
\setbeamertemplate{footline}{%
  \leavevmode\hbox to \paperwidth{\hspace{0.75cm}%
  {\color{gris}\tiny Francisco Castel \textperiodcentered{} Control y Sistemas \textperiodcentered{} Agitador magnético BLDC sensorless}%
  \hfill{\color{gris}\tiny\insertframenumber/\inserttotalframenumber}\hspace{0.75cm}}\vspace{0.28cm}}

\newcommand{\titulito}[1]{{\color{azul}\small\bfseries #1}\par\vspace{0.3em}}
\newcommand{\nro}[1]{\tikz[baseline=(n.base)]\node[circle, fill=acento, text=white,
  font=\tiny\bfseries, inner sep=0pt, minimum size=1.15em] (n) {#1};}
\newcommand{\rta}{{\color{acento}$\Rightarrow$}\hspace{0.3em}}

\pgfplotsset{
  estilo/.style={
    width=\linewidth, height=4.3cm,
    axis lines=left, axis line style={gris!70}, tick style={gris!70},
    tick label style={font=\scriptsize}, label style={font=\scriptsize},
    grid=major, grid style={grisclaro},
    /pgf/number format/.cd, use comma, 1000 sep={},
  },
  % Dibuja sólo las filas cuya columna #1 vale #2.
  discard if not/.style 2 args={
    x filter/.append code={%
      \edef\tempa{\thisrow{#1}}\edef\tempb{#2}%
      \ifx\tempa\tempb\else\def\pgfmathresult{inf}\fi}},
}

\begin{document}

% =====================================================================
% 1. Título
% Guion: el agitador tiene dos zonas donde el control "no ve": el rotor a
% baja velocidad y la píldora siempre. De eso trata el avance.
% =====================================================================
\begin{frame}[plain,c]
  \begin{columns}[c]
    \begin{column}{0.56\textwidth}
      {\color{acento}\scriptsize\bfseries CONTROL Y SISTEMAS \textperiodcentered{} TRABAJO FINAL}\par
      \vspace{0.7em}
      {\color{azul}\LARGE\bfseries Agitador magnético con\\[0.1em] motor BLDC sensorless\par}
      \vspace{0.8em}
      {\normalsize Dos no linealidades que limitan el control:\\ la zona ciega y el desenganche de la píldora\par}
      \vspace{1.6em}
      {\small\bfseries Francisco Castel}\par
      {\small Facultad de Ingeniería \textperiodcentered{} Universidad Nacional de Cuyo}\par
      \vspace{0.4em}
      {\small\color{gris}Consulta de avance \textperiodcentered{} 27 de noviembre de 2026}\par
    \end{column}
    \begin{column}{0.42\textwidth}
      \centering
      \begin{tikzpicture}[line cap=round, scale=1.12]
        % vaso con agua
        \fill[acento!12] (-1.42,1.93) rectangle (1.42,3.35);
        \draw[gris!80!black, line width=1pt] (-1.48,3.75) -- (-1.48,1.87) -- (1.48,1.87) -- (1.48,3.75);
        \node[font=\tiny, text=acento!70!black] at (0,3.12) {fluido};
        % píldora (polos opuestos al imán impulsor)
        \fill[azul!55] (-0.72,1.95) rectangle (0,2.15);
        \fill[rojo!60] (0,1.95) rectangle (0.72,2.15);
        \draw[gris!80!black] (-0.72,1.95) rectangle (0.72,2.15);
        \draw[-{Stealth[length=3.5pt]}, acento, thick]
          ($(0,2.05)+(200:0.92 and 0.2)$) arc[start angle=200, end angle=340, x radius=0.92, y radius=0.2];
        % base
        \fill[grisclaro] (-1.85,1.66) rectangle (1.85,1.81);
        \draw[gris] (-1.85,1.66) rectangle (1.85,1.81);
        % imán impulsor
        \fill[rojo!60] (-0.78,1.27) rectangle (0,1.47);
        \fill[azul!55] (0,1.27) rectangle (0.78,1.47);
        \draw[gris!80!black] (-0.78,1.27) rectangle (0.78,1.47);
        \draw[-{Stealth[length=3.5pt]}, azul, thick]
          ($(0,1.37)+(200:0.98 and 0.2)$) arc[start angle=200, end angle=340, x radius=0.98, y radius=0.2];
        % eje y motor
        \fill[gris!60] (-0.06,1.0) rectangle (0.06,1.27);
        \fill[grisclaro] (-0.85,0.1) rectangle (0.85,1.0);
        \draw[gris] (-0.85,0.1) rectangle (0.85,1.0);
        \node[font=\tiny, align=center] at (0,0.55) {motor BLDC\\(disco rígido)};
        % rótulos
        \draw[gris] (0.8,2.05) -- (1.9,2.45);
        \node[font=\tiny, anchor=west, align=left] at (1.9,2.45) {píldora: $\theta_2,\ \omega_2$\\{\color{rojo}no se mide}};
        \draw[gris] (0.85,1.37) -- (1.9,1.05);
        \node[font=\tiny, anchor=west, align=left] at (1.9,1.05) {imán: $\theta_1,\ \omega_1$\\{\color{acento!70!black}se mide (BEMF)}};
        \node[font=\tiny, anchor=west, text=azul] at (1.9,1.75) {$\delta=\theta_1-\theta_2$};
      \end{tikzpicture}
    \end{column}
  \end{columns}
\end{frame}

% =====================================================================
% 2. Avances
% Guion: qué anda hoy y qué preguntas del anteproyecto ya tienen respuesta.
% La medición nueva que cambia todo es la corriente de bus.
% =====================================================================
\begin{frame}{Qué está hecho y qué preguntas del anteproyecto ya tienen respuesta}
  \begin{columns}[T]
    \begin{column}{0.56\textwidth}
      \centering
      \begin{tikzpicture}[
          bloque/.style={draw=azul, fill=azul!5, rounded corners=2pt, align=center,
            font=\scriptsize, minimum height=0.85cm, inner sep=2.5pt},
          sensor/.style={draw=acento, fill=acento!7, rounded corners=2pt, align=center,
            font=\scriptsize, minimum height=0.65cm, inner sep=2.5pt},
          flecha/.style={-{Stealth[length=4pt]}, thick, gris!90!black},
          lbl/.style={font=\tiny, text=gris!90!black}]
        \node[bloque, minimum width=2.0cm] (mcu) at (0,0) {\textbf{STM32F103}\\PI 2 ms \textperiodcentered{} six-step\\barrido de carga};
        \node[bloque, minimum width=1.2cm] (puente) at (2.55,0) {\textbf{L298N}};
        \node[bloque, minimum width=1.4cm] (motor) at (4.45,0) {\textbf{Motor}\\BLDC};
        \node[bloque, minimum width=1.35cm, draw=rojo, fill=rojo!5] (pildora) at (6.35,0) {\textbf{Píldora}\\en el vaso};
        \node[sensor] (lm) at (2.55,-1.2) {LM339: cruces por cero};
        \node[sensor] (ina) at (3.0,1.2) {INA219: $V_{bus}$, $I_{bus}$ de la fuente};
        \node[bloque, minimum width=1.5cm, minimum height=0.65cm] (gui) at (-0.35,1.2) {\textbf{GUI} (PC)};
        \draw[flecha] (mcu) -- node[lbl, above] {PWM} node[lbl, below] {34 kHz} (puente);
        \draw[flecha] (puente) -- node[lbl, above] {3 fases} (motor);
        \draw[flecha, rojo, dashed] (motor) -- node[lbl, above, text=rojo] {imán} (pildora);
        \draw[flecha] (motor.south) |- (lm.east);
        \draw[flecha] (lm.west) -| (mcu.south);
        \draw[flecha] (ina.south -| puente.north) -- (puente.north);
        \draw[flecha, dashed] (ina.west) -| ([xshift=0.75cm]mcu.north);
        \draw[{Stealth[length=4pt]}-{Stealth[length=4pt]}, thick, gris!90!black]
          (gui.south) -- node[lbl, left] {USB} (gui.south |- mcu.north);
      \end{tikzpicture}

      \vspace{0.6em}
      \raggedright
      \titulito{Hecho}
      \scriptsize
      \begin{itemize}\setlength{\itemsep}{0.1em}
        \item Arranque SPWM a lazo abierto y paso a six-step por cruces por cero.
        \item PI de velocidad cada 2 ms, con rampa de consigna.
        \item Tensión y corriente de bus con el INA219.
        \item Firmware simplificado, sin HIL/PIL ni shunts: flash 79,6 \% $\to$ 76,2 \%.
        \item Barrido de carga en el firmware y página de curvas en la GUI.
        \item Curvas en aire y en agua: la píldora se suelta cerca de 2000 rpm.
      \end{itemize}
    \end{column}
    \begin{column}{0.42\textwidth}
      \titulito{Del anteproyecto a hoy}
      \footnotesize
      \textbf{¿Alcanza con medir la velocidad del motor?}\\
      \rta Para regular, sí. Para saber si la píldora gira, no: si se suelta, el PI sigue en consigna.

      \vspace{0.6em}
      \textbf{¿Corriente de fase o de bus?}\\
      \rta De bus: el INA219 alcanza para estimar la carga y ver el desenganche.

      \vspace{0.6em}
      \textbf{¿Conmutación actual o basada en corriente?}\\
      \rta Six-step por cruces por cero, sin sensado de fase.

      \vspace{0.6em}
      \textbf{¿Qué se valida con ensayos?}\\
      \rta Las curvas de carga, medidas en banco. De ellas salen los parámetros del modelo.
    \end{column}
  \end{columns}
\end{frame}

% =====================================================================
% 3. No linealidad 1: zona ciega
% Guion: la BEMF es proporcional a la velocidad; a baja velocidad el
% comparador no ve cruces. No es ruido: la medición desaparece. Por eso el
% arranque es a lazo abierto, y ahí el rotor puede perder el paso.
% Los valores de la rampa SPWM son los de hard_config.h.
% =====================================================================
\begin{frame}{No linealidad 1: a baja velocidad no hay medición}
  \begin{columns}[T]
    \begin{column}{0.47\textwidth}
      \vspace{0.6em}
      \begin{tikzpicture}[baseline=(current bounding box.north)]
        \begin{axis}[estilo, height=5.2cm, xmin=0, xmax=1500, ymin=0, ymax=1.3,
            xlabel={$n$ [rpm]}, ylabel={BEMF por fase, $e = k_e\,\omega_1$ [V]},
            xtick={0,1000,1500}, ytick={0,0.5,1},
            extra x ticks={350}, extra x tick labels={$n_{ciego}$},
            extra x tick style={tick label style={text=rojo}, grid=none}, clip=false]
          \fill[rojo!14] (axis cs:0,0) rectangle (axis cs:300,1.3);
          \fill[rojo!7] (axis cs:300,0) rectangle (axis cs:400,1.3);
          \fill[gris!12] (axis cs:0,0) rectangle (axis cs:1500,0.29);
          \draw[rojo, dashed, thick] (axis cs:0,0.29) -- (axis cs:1500,0.29)
            node[pos=0.74, above, font=\tiny, text=rojo] {umbral efectivo (LM339 + ruido PWM)};
          \addplot[azul, very thick, domain=0:1500, samples=2] {0.000827*x};
          \node[font=\tiny\bfseries, text=rojo, rotate=90] at (axis cs:150,0.85) {zona ciega};
          \node[font=\tiny, align=center, text=gris!80!black] at (axis cs:960,1.15) {cruces por cero válidos\\$\Rightarrow$ lazo cerrado};
          \node[font=\tiny, anchor=west, text=azul] at (axis cs:1060,0.7) {$k_e = K/2$, $K$ medido};
        \end{axis}
      \end{tikzpicture}
    \end{column}
    \begin{column}{0.51\textwidth}
      \small
      \begin{itemize}\setlength{\itemsep}{0.6em}
        \item La BEMF crece con la velocidad: $e = k_e\,\omega_1$. A baja velocidad queda
          tapada por el umbral y el ruido del comparador.
        \item Por debajo de $n_{ciego} \approx 300$--$400$ rpm no hay cruces por cero:
          \alert{la medición desaparece}. No es un error chico: es pérdida de observabilidad.
        \item Por eso el arranque es \textbf{a lazo abierto}: campo rotante SPWM
          (por defecto, $f_e$ de 2,1 a 9,3 Hz en 3 s). El rotor sigue al campo como un motor
          síncrono, $T \propto \sin(\theta_s - p\,\theta_1)$: con mucha carga o aceleración
          \alert{pierde el paso} y el firmware corta y reintenta.
        \item Por encima de $n_{ciego}$ la medición llega con un retardo $\tau \propto 1/n$:
          a baja velocidad el lazo tiene que ser más lento.
      \end{itemize}
    \end{column}
  \end{columns}
\end{frame}

% =====================================================================
% 4. No linealidad 2: acople magnético
% Guion: el acople es un resorte que transmite como máximo T_max. Si la carga
% lo supera no hay equilibrio y la píldora patina. Es el mismo seno que el
% arranque. El PI no lo ve; lo delata la corriente de bus (medido en agua).
% =====================================================================
\begin{frame}{No linealidad 2: el acople magnético transmite como máximo $T_{\max}$}
  \begin{columns}[T]
    \begin{column}{0.44\textwidth}
      \begin{tikzpicture}[baseline=(current bounding box.north)]
        \begin{axis}[estilo, height=4.1cm, xmin=0, xmax=3.35, ymin=0, ymax=1.32,
            xlabel={desfase $N_a\,\delta$}, ylabel={torque transmitido $T_c$},
            xtick={0,1.5708,3.1416}, xticklabels={$0$,$\pi/2$,$\pi$},
            ytick={0.55,1}, yticklabels={$T_{L,1}$,$T_{\max}$}, clip=false]
          \addplot[azul, very thick, domain=0:3.1416, samples=80] {sin(deg(x))};
          \draw[acento, thick, dashed] (axis cs:0,0.55) -- (axis cs:3.35,0.55);
          \draw[rojo, thick, dashed] (axis cs:0,1.15) -- (axis cs:3.35,1.15)
            node[pos=0.5, above, font=\tiny, text=rojo] {$T_{L,2} > T_{\max}$: no hay equilibrio, patina};
          \node[circle, fill=acento, inner sep=1.7pt] at (axis cs:0.5824,0.55) {};
          \node[circle, draw=acento, fill=white, inner sep=1.5pt] at (axis cs:2.5592,0.55) {};
          \node[font=\tiny, text=acento!70!black, below right] at (axis cs:0.6,0.53) {estable};
          \node[font=\tiny, text=gris, below left] at (axis cs:2.55,0.53) {inestable};
        \end{axis}
      \end{tikzpicture}

      \footnotesize
      \[
        J_2\,\dot\omega_2 = T_{\max}\sin(N_a\delta) - T_{L,tot}(\omega_2),
        \qquad \dot\delta = \omega_1 - \omega_2
      \]
      Hay equilibrio sólo si $T_{L,tot}(\omega) \le T_{\max}$. Cerca del límite, la rigidez
      del acople, $K_{eff} = N_a\sqrt{T_{\max}^2 - T_{L,tot}^2}$, tiende a cero.
    \end{column}
    \begin{column}{0.54\textwidth}
      \begin{tikzpicture}[baseline=(current bounding box.north)]
        \begin{axis}[estilo, height=4.1cm, xmin=900, xmax=3100, ymin=70, ymax=116,
            xlabel={$n$ [rpm]}, ylabel={$I_{bus}$ [mA]}]
          \fill[rojo!9] (axis cs:2000,70) rectangle (axis cs:2250,116);
          \addplot[gris, thick, mark=o, mark size=1.4pt] table[x=rpm, y=ibus] {datos/aire.dat};
          \addplot[azul, thick, mark=*, mark size=1.6pt, discard if not={acoplado}{1}] table[x=rpm, y=ibus] {datos/agua.dat};
          \addplot[rojo, thick, mark=*, mark size=1.6pt, discard if not={acoplado}{0}] table[x=rpm, y=ibus] {datos/agua.dat};
          \node[font=\tiny, text=azul, anchor=south] at (axis cs:1560,105) {agua, enganchada};
          \node[font=\tiny, text=gris!90!black, anchor=north west] at (axis cs:1250,78) {tara (vaso sin agua)};
          \node[font=\tiny, text=rojo, anchor=north] at (axis cs:2700,86) {agua, suelta};
          \node[font=\tiny, text=rojo, rotate=90] at (axis cs:2125,104) {se suelta};
        \end{axis}
      \end{tikzpicture}

      \footnotesize
      \begin{itemize}\setlength{\itemsep}{0.25em}
        \item \textbf{Mismo seno que en el arranque}: son dos acoples síncronos en serie
          (campo $\to$ rotor e imán $\to$ píldora) y en ninguno se mide el ángulo.
        \item \textbf{El PI no lo ve en régimen}: regula $\omega_1$ y vuelve a la consigna.
          Lo delatan $I_{bus}$, que cae por debajo de la tara, y la dispersión de velocidad
          (de 5--9 a 24--38 rpm).
      \end{itemize}
    \end{column}
  \end{columns}
\end{frame}

% =====================================================================
% 5. Propuesta: curvas de carga
% Guion: método en cuatro pasos y lo que ya muestra la curva de agua. Cerrar
% con los ensayos y cuáles se hacen antes de la mesa.
% =====================================================================
\begin{frame}{Propuesta: curvas de carga por fluido y margen al desenganche}
  \begin{columns}[T]
    \begin{column}{0.53\textwidth}
      \titulito{Metodología}
      \footnotesize
      \begin{enumerate}\setlength{\itemsep}{0.45em}
        \item \textbf{Barrido en escalera} (firmware): media y dispersión de $n$, $d$,
          $V_{bus}$ e $I_{bus}$ en cada escalón.
        \item \textbf{Tara} con el vaso sin líquido. Como $I_{bus} \approx I_0 + d\,I_{fase}$
          (tara: $50 + 106\,d$ mA), restarla a igual velocidad cancela el consumo propio $I_0$.
        \item \textbf{Carga del fluido}: $\Delta T = K_t\,(I_{bus} - I_{bus,tara})/d$, con
          $K_t = 0{,}0158$ N$\cdot$m/A, la pendiente de $d\,V_{bus}$ contra $\omega$ en la tara.
        \item \textbf{Límites}: ajustar $T_L(\omega)$ y operar con margen $m$:
          $T_{L,tot}(\omega) \le m\,T_{\max}$ y rampa
          $\alpha \le \bigl(m\,T_{\max} - T_{L,tot}(\omega)\bigr)/J_{2,eff}$.
      \end{enumerate}
    \end{column}
    \begin{column}{0.45\textwidth}
      \begin{tikzpicture}[baseline=(current bounding box.north)]
        \begin{axis}[estilo, height=3.7cm, xmin=1100, xmax=2900, ymin=-0.75, ymax=1.15,
            xlabel={$n$ [rpm]}, ylabel={$\Delta T$ [mN$\cdot$m]}, ytick={-0.5,0,0.5,1}]
          \fill[rojo!9] (axis cs:2000,-0.75) rectangle (axis cs:2250,1.15);
          \draw[gris, thick] (axis cs:1100,0) -- (axis cs:2900,0)
            node[pos=0.93, above, font=\tiny, text=gris] {tara};
          \addplot[azul, thick, mark=*, mark size=1.6pt, discard if not={acoplado}{1}] table[x=rpm, y=dT] {datos/agua.dat};
          \addplot[rojo, thick, mark=*, mark size=1.6pt, discard if not={acoplado}{0}] table[x=rpm, y=dT] {datos/agua.dat};
          \node[font=\tiny, text=azul, anchor=north] at (axis cs:1640,0.66) {enganchada: $+0{,}75$ a $+0{,}87$};
          \node[font=\tiny, text=rojo, anchor=south] at (axis cs:2580,-0.3) {suelta};
        \end{axis}
      \end{tikzpicture}

      \scriptsize
      Ya se ve algo: $\Delta T$ crece mucho menos que $\omega^2$. Con aire y agua solos no se
      separa la fricción de la píldora del arrastre del fluido: falta la tara sin píldora.
    \end{column}
  \end{columns}

  \vspace{0.6em}
  {\color{grisclaro}\hrule height 0.6pt}
  \vspace{0.45em}
  \scriptsize
  \textbf{\color{azul}Ensayos}\enspace
  \nro{1} repetibilidad\enspace
  \nro{2} tara sin píldora: $T_{\max} = \Delta T(n_{dec}) + T_{f2}$\enspace
  \nro{3} desenganche fino, 1800--2300 rpm cada 50\\[0.35em]
  \hspace*{3.05em}\nro{4} rampa contra desenganche\enspace
  \nro{5} volumen de líquido\enspace
  \nro{6} glicerina (viscosidad)\hfill
  {\color{acento}\bfseries Antes de la mesa: 1, 2 y 3}
\end{frame}

% =====================================================================
% 6. Utilidad del modelo
% Guion: el modelo no es una simulación con parámetros inventados. Son pocas
% ecuaciones, cada parámetro sale de una curva, y sirven para decidir: qué
% leer en la corriente, hasta qué velocidad y con qué rampa operar, dónde
% vale el PI y qué ensayo hace falta.
% =====================================================================
\begin{frame}{¿Para qué sirve el modelo? Para pasar de curvas a límites de operación}
  \centering
  \begin{tikzpicture}[x=0.0041cm, y=1cm]
    \fill[rojo!14] (0,0) rectangle (350,0.5);
    \fill[verde!16] (350,0) rectangle (2000,0.5);
    \shade[left color=verde!16, right color=rojo!14] (2000,0) rectangle (2250,0.5);
    \fill[rojo!14] (2250,0) rectangle (3300,0.5);
    \draw[gris] (0,0) rectangle (3300,0.5);
    \node[font=\tiny\bfseries, text=rojo] at (175,0.25) {zona ciega};
    \node[font=\scriptsize\bfseries, text=verde!70!black] at (1175,0.25) {modelo lineal y PI válidos};
    \node[font=\scriptsize\bfseries, text=rojo] at (2775,0.25) {la píldora patina};
    \foreach \x in {0,1000,2000,3000} {\draw[gris] (\x,0) -- (\x,-0.07) node[below, font=\tiny, text=gris] {\x};}
    \node[font=\tiny, text=gris, anchor=west] at (3320,-0.2) {rpm};
    \node[font=\tiny, anchor=south west] at (0,0.52) {$n_{ciego}\approx 300$--$400$ rpm (sensor)};
    \node[font=\tiny, anchor=south] at (2125,0.52) {$n_{dec}$ en agua $\approx 2000$--$2250$ rpm (acople y fluido)};
    \draw[-{Stealth[length=3.5pt]}, rojo, thick] (2050,-0.42) -- (1650,-0.42)
      node[left, font=\tiny, text=rojo] {con más carga, $n_{dec}$ baja};
  \end{tikzpicture}

  \vspace{0.3em}
  \raggedright
  \begin{columns}[T]
    \begin{column}{0.49\textwidth}
      \titulito{Modelo simple, con parámetros medidos}
      \footnotesize
      \begin{itemize}\setlength{\itemsep}{0.3em}
        \item Eléctrico casi estático ($L/R \ll \tau_{mec}$):
          $d\,V_{bus} \approx K\,\omega_1 + R_{eq}\,I + V_{L298}$. $K$ sale de la tara.
        \item Mecánico, con la píldora enganchada: $G_1(s) = \dfrac{K_{dc}}{\tau_{mec}\,s + 1}$
          para el PI. Vale si el lazo es bastante más lento que la resonancia del acople.
        \item Acople $T_{\max}\sin(N_a\delta)$: se usa como \textbf{límite}; no hace falta simularlo en detalle.
        \item Carga $T_L(\omega)$: la dan las curvas, no coeficientes genéricos.
      \end{itemize}
    \end{column}
    \begin{column}{0.49\textwidth}
      \titulito{Qué aporta}
      \footnotesize
      \begin{enumerate}\setlength{\itemsep}{0.3em}
        \item \textbf{Leer la medición}: $I_{bus}$ se convierte en torque y la tara cancela $I_0$.
        \item \textbf{Predecir} $n_{dec}$ y la rampa segura antes de que la píldora se suelte.
        \item \textbf{Acotar el diseño}: el PI vale para $n_{ciego} < n < n_{dec}$.
        \item \textbf{Elegir ensayos}: aire y agua no identifican la viscosidad (el agua está en
          régimen inercial, $Re \sim 10^4$). Hace falta un fluido viscoso de referencia.
      \end{enumerate}
    \end{column}
  \end{columns}

  \vspace{0.3em}
  \centering
  \colorbox{acento!10}{\parbox{0.94\textwidth}{\centering\small
    Las curvas le dan los parámetros al modelo; el modelo dice qué consigna y qué rampa son seguras.}}
\end{frame}

\end{document}
